\section{ISRF channel: packet-hop transport}
\label{sec:packet}

This section describes the packet-hop transport of the two ISRF bands, which
\code{GEARFeedback:ISRF\_transport: packet\_hop} selects in place of the M1
system of Section~\ref{sec:isrf-eqdisc}. It needs \code{GEARFeedback:ISRF\_propagation}
and the HEALPix C library. It runs in static and in cosmological boxes. Positions,
offsets, smoothing lengths and path lengths are comoving, and a physical length
is $a$ times the comoving one. Dust absorption, the redshift of the light, the step
lengths and the field that the gas sees are physical.

A packet is a quantity of band energy with a photon number and a direction.
A star hands its light to the gas around it. Each gas particle holds the light
it received, and hands it on to its downstream neighbours once per own step.
The transport has no propagation speed: one hop is one step of the sending
particle, whatever the distance. The sections below give the packet
(Section~\ref{sec:packet-state}), the hop (\ref{sec:packet-hop}), the merge
(\ref{sec:packet-merge}), the emission by a star (\ref{sec:packet-emission}),
the redshift and the band-edge transfer (\ref{sec:packet-transfer}), the dust
absorption (\ref{sec:packet-absorption}), the treatment of unequal
time bins (\ref{sec:packet-reservoir}), the field that the gas sees
(\ref{sec:packet-field}), the frame of reference of the light
(\ref{sec:packet-frame}) and the ledgers (\ref{sec:packet-ledger}).

\subsection{Symbols and conventions}
\label{sec:packet-symbols}

Bands are $b\in\{{\rm PE},{\rm LW}\}$. The band lower edges are
$E_{{\rm lo},{\rm PE}}=6.0$\,eV and $E_{{\rm lo},{\rm LW}}=11.2$\,eV, and the
upper edges $E_{{\rm hi},b}$ are $11.2$\,eV (PE) and $13.6$\,eV (LW). Slots are numbered by $k$,
sending particles by $i$, receiving particles by $j$. Lengths are comoving unless
marked physical.

\begin{center}
\begin{tabular}{p{1.9cm}p{3.4cm}p{8.6cm}}
\hline
symbol & units & meaning \\
\hline
$E_{k,b}$ & energy & band energy held in slot $k$ of a particle \\
$N_{k,b}$ & energy & photon number of slot $k$, stored as $N E_{{\rm lo},b}$ (number of photons times the lower edge energy) \\
$\bar E_b$ & energy & mean photon energy of the band, $E_{{\rm lo},b}\,E_{k,b}/N_{k,b}$ \\
$f_b$ & -- & photon number per unit energy of a star, $E_{{\rm lo},b}/\bar E_b$ \\
$\boldsymbol{n}_k$ & -- & unit direction of travel of slot $k$ \\
$d_k$ & length & distance from the origin of slot $k$ to the particle \\
$\boldsymbol{o}$ & length & offset from a particle to the origin of a packet \\
$\boldsymbol{\Delta x}_{ij}$ & length & vector from sender $i$ to receiver $j$ \\
$h_i$ & length & smoothing length of sender $i$; kernel support $\gamma h_i$ \\
$W,\gamma$ & length$^{-3}$, -- & hydro kernel of the build and its support in units of $h$ \\
$p$ & -- & power of the hop cone, \code{ISRF\_packet\_cone\_power} (default 4) \\
$\mu_{ij,k}$ & -- & cosine between $\boldsymbol{\Delta x}_{ij}$ and $\boldsymbol{n}_k$ \\
$V_j$ & length$^3$ & comoving volume $m_j/\rho_j$ of receiver $j$ \\
$w_{ij,k}$ & -- & hop weight of slot $k$ towards $j$ \\
$S_{i,k}$ & -- & sum of the hop weights of slot $k$ over the neighbours of $i$ \\
$\kappa_b$ & 1/length & linear dust absorption rate of band $b$ at the receiver, physical \\
$\lambda_E(b)$, $\lambda_N$ & -- & band-edge weight of the energy of band $b$, and of the LW photon number (Section~\ref{sec:packet-transfer}) \\
$H$, $\zeta$ & 1/time, -- & Hubble rate, and redshift depth of a share, $\zeta=H\,\ell_{\rm phys}/c$ \\
$\ell_{\rm phys}$ & length & physical path length of a share, $\ell_{\rm phys}=a\,\Delta\ell$ \\
$\nu$ & -- & rate per unit $\zeta$ at which the PE photon number crosses $6$\,eV \\
$R$ & -- & rotation matrix of the slot frame \\
$E_{{\rm hi},b}$ & energy & upper edge of band $b$ \\
$L_b$ & energy/time & band luminosity of a star, physical \\
$\Delta t_\star$ & time & feedback step of the star, proper time \\
$m_j$, $\rho_j$ & mass, mass/length$^3$ & mass and (comoving) density of gas particle $j$ \\
$a$ & -- & scale factor, 1 in a static box \\
$f_{{\rm rel},k}$ & -- & fraction of the reservoir of slot $k$ released at a merge \\
$B_{k,b}$ & energy & reservoir of slot $k$ (light from sources with a longer step) \\
$S_{\rm floor}$ & length$^3$ & cell floor of a star's emission \\
$E_{\rm fl}$ & energy & light in flight \\
$Y$ & energy & cumulative energy removed by the expansion from the shares booked into a gas particle, net of the gain from the other band (Section~\ref{sec:packet-transfer}) \\
$I_{\rm inh}$ & energy & inherited ledger of a star or a sink \\
\hline
\end{tabular}
\end{center}

\subsection{The packet and its reception slots}
\label{sec:packet-state}

A gas particle holds one packet per reception slot. The slots are the cells of a
partition of the sphere of directions into $K$ cells of equal solid angle. The
cells are convex cones fixed in a rotated frame, so the mean of any set of
directions of one cell lies in that cell. Each slot holds, for both bands, the
energy $E_{k,b}$ and the photon number $N_{k,b}$, and for both bands together
one direction $\boldsymbol{n}_k$ and one distance $d_k$ to its origin. The mean photon energy
of a band follows from the pair,
\begin{equation}
  \bar E_b = E_{{\rm lo},b}\,\frac{E_{k,b}}{N_{k,b}}.
  \label{eq:packet-meanE}
\end{equation}
A star sets the photon number of its emission from the mean photon energy
of its spectrum in the band, clamped to the band,
\begin{equation}
  N = f_b\,E, \qquad f_b=\frac{E_{{\rm lo},b}}{\bar E_b}, \qquad
  E_{{\rm lo},b}\le \bar E_b\le E_{{\rm hi},b}.
  \label{eq:packet-photonfrac}
\end{equation}

\paragraph{The slot function.} Let $\boldsymbol{y}=R\,\boldsymbol{x}$ be a
direction $\boldsymbol{x}$ in the frame, $R$ a rotation matrix. The slot of
a direction depends on the number $K$ of slots, a compile-time option.
\begin{align}
  K=8:&\quad k=\sum_{m=0}^{2} 2^m\,[\,y_m\ge 0\,], \label{eq:packet-slot8}\\
  K=12:&\quad k=4m+2\,[\,y_{i}\ge0\,]+[\,y_{j}\ge0\,], \quad
  m=\arg\min_{m'} |y_{m'}|, \label{eq:packet-slot12}\\
  K=16:&\quad k=\sum_{m=0}^{2} 2^m\,[\,y_m\ge 0\,]+8\,[\,|y_1|>|y_0|\,],
  \label{eq:packet-slot16}
\end{align}
with $[\cdot]$ equal to 1 if the condition holds and 0 otherwise. In
Eq.~\eqref{eq:packet-slot12}, $m$ is the axis of the smallest $|y|$, the lowest axis on a tie, and $(i,j)$ are the two
other axes in axis order (here $i$ and $j$ are axis indices): $(1,2)$ for $m=0$, $(0,2)$ for $m=1$ and $(0,1)$ for
$m=2$. For $K=8$ the cells are the octants of the frame. For $K=12$
they are the faces of a rhombic dodecahedron. For $K=16$ they are the octants
halved by the plane $|y_0|=|y_1|$; a tie $|y_0|=|y_1|$ goes to the lower half. A direction on a border goes to the cell
with the positive coordinate, whatever the sign of a zero. A
direction in the same octant as another differs from it by at most
$90^\circ$, so for $K=8$ light from opposite sides is never put in one slot.

\paragraph{The frame.} The matrix is
\begin{equation}
  R = R_z(\theta_z)\,R_y(\theta_y)\,R_x(\theta_x),
  \label{eq:packet-frame}
\end{equation}
with $\theta_x,\theta_y,\theta_z$ the angles of
\code{ISRF\_packet\_octant\_frame\_angles\_deg}, default $(24^\circ,36^\circ,66^\circ)$.
The default keeps every cell border at least $10^\circ$ from the six box axes for $K=12$ ($19^\circ$ for $K=8$, $6^\circ$ for $K=16$). Box diagonals can lie on a border: the $(1,1,1)$ directions are $0.01^\circ$ from one at $K=12$. A source or beam exactly along a box diagonal then gets an arbitrary but deterministic slot; rotate the frame with \code{ISRF\_packet\_octant\_frame\_angles\_deg} for such a fixture. A positive seed offset
\code{ISRF\_packet\_octant\_frame\_seed\_offset} multiplies $R$ on the left by
a rotation drawn uniformly from that seed \citep{Shoemake1992}, so that repeats of a
run use independent frames. The frame is the same for every particle.

\paragraph{Implementation.} All locations are in the repository root. Line
numbers are from branch \code{darwin/feat/isrf\_cone}; where lines move, find
the construct by name. The slot count is
\code{RADIATION\_PACKET\_SLOTS}, set by \code{--with-isrf-packet-slots=<N>}
(\code{configure.ac:3048-3056}, default 12; \code{N} other than 8, 12 or 16
is an error; the default and the check are repeated at
\code{src/feedback/GEAR\_thermal/feedback\_struct.h:200-207}). The slot
struct is \code{feedback\_isrf\_packet\_slot\_data}
(\code{feedback\_struct.h:310-338}); the packet is
\code{feedback\_isrf\_packet\_data} (\code{feedback\_struct.h:344-372}).
Eq.~\eqref{eq:packet-meanE} is the definition in the comment at
\code{feedback\_struct.h:237-239}.
Eq.~\eqref{eq:packet-photonfrac} is
\code{radiation\_set\_star\_packet\_photon\_fraction}
(\code{src/feedback/GEAR/radiation.h:389-401}, the clamp at :398-399, the ratio at :400). The transport stops at
start-up if a yields table lacks \code{MeanPhotonEnergyPE} or
\code{Integrated\_MeanPhotonEnergyPE}
(\code{src/feedback/GEAR\_thermal/feedback\_properties.h:1102-1110}).
Eq.~\eqref{eq:packet-slot8} is \code{radiation\_packet\_octant}
(\code{src/feedback/GEAR/radiation\_packet.c:220-231}), and
Eqs.~\eqref{eq:packet-slot8}--\eqref{eq:packet-slot16} are
\code{radiation\_packet\_slot} (\code{radiation\_packet.c:251-277}: $K=8$ at :254-255,
$K=16$ at :262-267, $K=12$ at :268-275). Eq.~\eqref{eq:packet-frame} is
\code{radiation\_packet\_build\_octant\_frame}
(\code{radiation\_packet.c:330-386}: the product at :346-356, the extra rotation at
:358-382). The frame is copied onto each particle by
\code{radiation\_set\_part\_transport} (\code{radiation\_isrf.c:152-162}).

\subsection{The hop}
\label{sec:packet-hop}

Once per own step, every active sender $i$ hands the energy and the photon
number held in each non-empty slot $k$ to its neighbours. Slot $k$ sends to
neighbour $j$ the fraction
\begin{equation}
  \phi_{ij,k} = \frac{w_{ij,k}}{S_{i,k}}, \qquad
  S_{i,k}=\sum_{j'} w_{ij',k}, \qquad
  w_{ij,k}=V_j\,W(r_{ij},h_i)\,\mu_{ij,k}^{\,p},
  \label{eq:packet-share}
\end{equation}
where $r_{ij}=|\boldsymbol{\Delta x}_{ij}|$ and
$\mu_{ij,k}=\boldsymbol{\Delta x}_{ij}\!\cdot\!\boldsymbol{n}_k/r_{ij}$. The sum
$S_{i,k}$ is the SPH quadrature of $\int W\mu^p\,\mathrm{d}V$, so across a density
gradient the shares follow the solid angle and not the number of neighbours on
each side. The density is the previous-step snapshot read by the gradient and the
force loop alike; the one out-of-line function divides $m_j$ by $\rho_j$ for both loops. Its expansion factor is common to all receivers and cancels in
$\phi_{ij,k}$. The weight $w_{ij,k}$ is zero if $r_{ij}=0$, if $r_{ij}\ge\gamma h_i$, or if
$\mu_{ij,k}\le\mu_{\rm min}=10^{-3}$. The sum runs over all neighbours with
non-zero weight. The sender's weights are equal in the gradient loop that
forms $S_{i,k}$ and in the force loop that books the shares, because one
function computes both, so $\sum_j\phi_{ij,k}=1$ up to float rounding. The
hop hands on the whole of each slot: nothing is kept by the sender.

The reach of a hop is the kernel support of the sender: a packet hops only to
neighbours with $r_{ij}<\gamma h_i$. The reach is a geometric bound and not a
light travel distance. The step of the sender does not limit it, and no
timestep bound from the speed of light applies (Section~\ref{sec:packet-field}).

For the Wendland-C2 kernel of the canonical build (three dimensions),
\begin{equation}
  W(r,h) = \frac{21}{2\pi\,\gamma^{3}h^{3}}\,(1-q)^4\,(1+4q), \qquad
  q=\frac{r}{\gamma h},\quad \gamma=1.936492,
  \label{eq:packet-kernel}
\end{equation}
which is the polynomial $4q^5-15q^4+20q^3-10q^2+1$, equal to $(1-q)^4(1+4q)$
\citep{Wendland1995}. Other kernels give their own $W$. The normalisation and the factor
$h_i^{-3}$ are the same for every neighbour of $i$, so they cancel in
$\phi_{ij,k}$: only the shape of $W$, the volume $V_j$ and the factor $\mu^p$ matter. A larger $p$
gives a narrower cone.

The share is placed in the slot of the receiver that holds the direction from
the packet's origin to the receiver. Seen from the receiver, the origin lies at
\begin{equation}
  \boldsymbol{o}_{j} = -\boldsymbol{\Delta x}_{ij} - d_{i,k}\,\boldsymbol{n}_{i,k},
  \label{eq:packet-offset}
\end{equation}
because the origin is $d_{i,k}$ behind the sender along its direction. The slot of the
receiver is the slot function of Eqs.~\eqref{eq:packet-slot8}--\eqref{eq:packet-slot16}
applied to $-\boldsymbol{o}_j$. The slot is
chosen from this geometry and not from the slot index of the sender. The share
carries the energy $\phi_{ij,k}E_{i,k,b}$ and the photon number $\phi_{ij,k}N_{i,k,b}$
of each band $b$, changed by the redshift and the band-edge transfer of Section~\ref{sec:packet-transfer} and reduced by the absorption of Section~\ref{sec:packet-absorption}.

The gradient loop also forms the mean advance of the slot along its direction,
\begin{equation}
  l_{i,k} = \frac{1}{S_{i,k}}\sum_{j} w_{ij,k}\,
  \boldsymbol{\Delta x}_{ij}\!\cdot\!\boldsymbol{n}_{i,k},
  \label{eq:packet-advance}
\end{equation}
which Section~\ref{sec:packet-field} uses. If a non-empty slot has
$S_{i,k}=0$, it has no downstream neighbour, and its light is booked as escaped
(Section~\ref{sec:packet-ledger}).

\paragraph{Implementation.} The weight, Eq.~\eqref{eq:packet-share} third
formula, is \code{radiation\_packet\_transmission\_weight}
(\code{radiation\_packet.c:196-212}): the zero cases at :202 and :206, $\mu$ at :205,
$W$ at :208-210, the volume $V_j=m_j/\rho_j$ and the factor $\mu^p$ at :211. The cutoff
$\mu_{\rm min}$ is \code{RADIATION\_PACKET\_MU\_MIN}
(\code{src/feedback/GEAR/radiation\_packet.h:53}). The function is out of line and divides
$m_j$ by $\rho_j$ inside, so both loops use the same compiled arithmetic. The coefficients of
Eq.~\eqref{eq:packet-kernel} are \code{kernel\_coeffs}, \code{kernel\_constant}
and \code{kernel\_gamma} in \code{src/kernel\_hydro.h:137-146}, evaluated by
\code{kernel\_eval} (\code{kernel\_hydro.h:297-318}). The gradient loop sums
$S_{i,k}$ and the numerator of Eq.~\eqref{eq:packet-advance} in
\code{radiation\_packet\_gradient\_accumulate}
(\code{radiation\_packet.h:155-175}), called from
\code{runner\_iact\_isrf\_gradient} and its non-symmetric twin
(\code{src/feedback/GEAR/radiation\_propagation\_iact.h:510-511}, :573). The division by
$S_{i,k}$ in Eq.~\eqref{eq:packet-advance} is at
\code{radiation\_packet.c:775}. The share $\phi_{ij,k}$ is
\code{fraction} in \code{radiation\_packet\_book\_share}
(\code{radiation\_packet.h:328}), and Eq.~\eqref{eq:packet-offset} is its
\code{offset} (:331-333). The receiving slot is chosen in
\code{radiation\_packet\_deliver} from the vector \code{from\_origin}
(\code{radiation\_packet.h:218}). The force loop calls
\code{radiation\_packet\_book\_share} from
\code{runner\_iact\_isrf\_packet\_hop} and its non-symmetric twin
(\code{radiation\_propagation\_iact.h:1043-1044}, :1075). The vector
\code{dx\_ij} passed to it points from sender to receiver. The sender's held
energy is set to zero at the end of the force loop
(\code{radiation\_packet\_end\_force}, \code{radiation\_packet.c:817-820}). The
escape of a slot with $S_{i,k}=0$ is at \code{radiation\_packet.c:785-792}. The parameter $p$ is
copied to each particle at \code{radiation\_isrf.c:158}, and
\code{ISRF\_packet\_cone\_power} is parsed at
\code{src/feedback/GEAR\_thermal/feedback\_properties.h:1037-1039}.

\subsection{The merge}
\label{sec:packet-merge}

At each active step of the particle, when its previous step has ended, in each slot, the energy and the photon
number that have arrived since the last active step become the held packet.
Let $I_{k,b}$ be the energy that arrived in the inbox of slot $k$, and
$f_{{\rm rel},k} B_{k,b}$ the part of the reservoir of Section~\ref{sec:packet-reservoir}
that is due. Then
\begin{equation}
  E_{k,b} \leftarrow E_{k,b} + I_{k,b} + f_{{\rm rel},k} B_{k,b},
  \label{eq:packet-merge-E}
\end{equation}
and the photon number $N_{k,b}$ is updated in the same way. Each arriving share
brought an offset $\boldsymbol{o}_s$ (Eq.~\eqref{eq:packet-offset}, or the star-to-gas
vector of Section~\ref{sec:packet-emission}) and a transmitted energy
$e_s=\sum_b e_{s,b}$ summed over both bands. The slot takes the
energy-weighted mean origin of its arrivals,
\begin{equation}
  \boldsymbol{o}_k = \frac{\sum_s e_s\,\boldsymbol{o}_s}{\sum_s e_s},
  \qquad d_k=|\boldsymbol{o}_k|, \qquad
  \boldsymbol{n}_k = -\frac{\boldsymbol{o}_k}{d_k},
  \label{eq:packet-merge-origin}
\end{equation}
for $d_k>0$. For $d_k=0$ the direction is $\boldsymbol{n}_k=(0,0,1)$. One origin and
one direction serve both bands. A slot that received no energy keeps its origin and direction. The energy held
before the merge is not part of the mean: the slot is empty at the merge,
because the force loop sent its content.

Every share was placed in its slot by its own geometry, so slots never merge
with each other. Within a slot, the merge conserves the energy and the photon
number of each band exactly (they are sums). It does not conserve the
directions: the individual directions of the arriving shares are replaced by
one mean direction, and the angular spread of the slot is lost. Because
$|\sum e_s\boldsymbol{o}_s|\le\sum e_s|\boldsymbol{o}_s|$, the merged distance $d_k$ is not
larger than the energy-weighted mean of the distances of the shares. The redshift
and the dust absorption of the later hops are computed from the merged origin
(Sections~\ref{sec:packet-transfer} and \ref{sec:packet-absorption}), not from the sum of the separate paths.
The vector $-\boldsymbol{o}_k$ is a positive combination of the vectors $-\boldsymbol{o}_s$, which all lie
in one convex cone, so it lies in that cone, and the direction $\boldsymbol{n}_k$ is in the cell
of its slot, apart from the case $d_k=0$.

\paragraph{Implementation.} The merge of one slot is
\code{radiation\_packet\_merge\_slot} (\code{radiation\_packet.c:600-692}), called
for every slot by \code{radiation\_packet\_merge} (:709-723), which
\code{radiation\_snapshot\_part\_propagation} calls for active particles
(\code{radiation\_isrf.c:257-266}). Eq.~\eqref{eq:packet-merge-E} is at
\code{radiation\_packet.c:644-647}. The sums of Eq.~\eqref{eq:packet-merge-origin} are
formed in \code{radiation\_packet\_deliver} as \code{inbox\_origin\_offset}
(\code{radiation\_packet.h:280-282}, weight \code{delivered\_energy}, summed
over the bands at :271); the division by $\sum_s e_s$ is at
\code{radiation\_packet.c:669}, $d_k$ at :670-672, $\boldsymbol{n}_k$ at :674, and the
fallback direction at :684-691. The merge returns before any change of origin
when the slot received no energy (:667). The fields of
Eq.~\eqref{eq:packet-merge-origin} are \code{origin\_distance},
\code{direction} and \code{inbox\_origin\_offset} in
\code{feedback\_isrf\_packet\_slot\_data} (\code{feedback\_struct.h:315-323}).

\subsection{Emission by a star}
\label{sec:packet-emission}

A star gives its light of one feedback step, $L_b\Delta t_\star$ per band, to the
gas particles around it. These particles are the receivers: gas with
$q<q_{\rm e}$, where $q=r/(\gamma h_\star)$ is the distance to the star over the kernel
support radius, $q_{\rm e}=0.95$, and a positive density. The star
does not use the M1 injection. The receivers are sorted into $n_c=12$ emission cells, the
pixels of a HEALPix map with \code{nside}$=1$. These cells have equal solid angle
$4\pi/12$. Their frame is a uniform random rotation $Q$ that depends on the ID of the star and on
the start $t_{\rm b}$ of its step, so the cell pattern changes at each step. With three
uniform numbers $u_1,u_2,u_3\in[0,1)$ drawn from a counter-based generator keyed by
the star ID and by $3t_{\rm b}+0$, $+1$, $+2$, the unit quaternion
$(q_x,q_y,q_z,q_w)$ is \citep{Shoemake1992}
\begin{align}
  q_x&=\sqrt{1-u_1}\,\sin 2\pi u_2, & q_y&=\sqrt{1-u_1}\,\cos 2\pi u_2,\nonumber\\
  q_z&=\sqrt{u_1}\,\sin 2\pi u_3, & q_w&=\sqrt{u_1}\,\cos 2\pi u_3,
  \label{eq:packet-quat}
\end{align}
and the cell $c$ of a gas particle is the pixel \citep{Gorski2005} that contains $Q\,\hat{\boldsymbol{d}}$,
with $\hat{\boldsymbol{d}}$ the unit vector from the star to the gas.

Each receiver $j$ in cell $c$ has the weight
\begin{equation}
  w_j = \frac{m_j}{\rho_j}\Big(1-\frac{q_j}{q_{\rm e}}\Big),
  \label{eq:packet-tent}
\end{equation}
which is the comoving volume of the particle times a linear (tent) function of
$q$ that falls to zero at $q_{\rm e}$. The weight is zero beyond $q_{\rm e}$, so
a pair near the edge of the support cannot change a share. Let $S_c=\sum_{j\in c}w_j$
be the weight sum of the cell, and
\begin{equation}
  S_{\rm floor} = f_0\,\frac{\pi}{3}\,\frac{(q_{\rm e}\gamma h_\star)^3}{n_c},
  \qquad f_0=\tfrac12,
  \label{eq:packet-floor}
\end{equation}
the cell floor. $S_{\rm floor}$ is $f_0$ times the value of $S_c$ for a uniform gas, because
$\int_0^{r_{\max}}4\pi r^2(1-r/r_{\max})\,dr=\pi r_{\max}^3/3$ with $r_{\max}=q_{\rm e}\gamma h_\star$.
Each cell holds $1/n_c$ of the star's light. The part of it that no receiver takes
escapes, with the escaping fraction of cell $c$
\begin{equation}
  \varepsilon_c =
  \begin{cases}
    0, & S_c\ge S_{\rm floor},\\
    (1-S_c/S_{\rm floor})/n_c, & 0<S_c<S_{\rm floor},\\
    1/n_c, & S_c=0.
  \end{cases}
  \label{eq:packet-escape}
\end{equation}
The delivered fraction $D=1-\sum_c\varepsilon_c$ is shared by all receivers of the star
in proportion to their weights: receiver $j$ gets the fraction of the star's light
\begin{equation}
  \varphi_j = D\,\frac{w_j}{S_{\rm tot}}, \qquad S_{\rm tot}=\sum_c S_c,
  \label{eq:packet-emit-share}
\end{equation}
so $\sum_j\varphi_j=D$. A cell with fewer emitters than its volume needs does not raise
the share of its receivers, because the shares do not depend on the cell of the receiver. A cell
with a real hole ($S_c<S_{\rm floor}$) still lets its part escape. A star with no receiver loses all its light, $D=0$.
The energy of receiver $j$ in
band $b$ is $L_b\Delta t_\star\,\varphi_j$ and its photon number is
$f_b$ times that energy (Eq.~\eqref{eq:packet-photonfrac}). The share arrives
with the offset $\boldsymbol{o}=\boldsymbol{x}_\star-\boldsymbol{x}_j$ of the origin at
the star, and it is absorbed along the distance $r=|\boldsymbol{o}|$
(Section~\ref{sec:packet-absorption}). The weight sums $S_c$ come from the star's
density loop and the shares from its feedback loop, which use the same function.

\paragraph{Implementation.} The star's density loop calls
\code{radiation\_packet\_star\_density} (called at
\code{src/feedback/GEAR\_thermal/feedback\_iact.h:75}), which adds the weight
to \code{cone\_weight[c]} (\code{radiation\_packet.h:453}). The feedback loop
calls \code{radiation\_packet\_star\_emit} (called at
\code{src/feedback/GEAR/radiation\_iact.h:275-277} in place of the M1 branch,
:278 on). Both loops take $c$, $q$ and $w_j$ from
\code{radiation\_packet\_star\_gas\_weight} (\code{radiation\_packet.c:551-576}):
$q$ at :568-570, the receiver test at :572-574, the wrapped star-gas vector at
:558-567. The cell is \code{radiation\_packet\_emission\_cone}
(:490-529): the numbers $u_1,u_2,u_3$ at :496-502, Eq.~\eqref{eq:packet-quat} at
:504-510, the rotation of $\hat{\boldsymbol{d}}$ at :512-520, the pixel at :522-524. The
number of cells is \code{RADIATION\_PACKET\_EMISSION\_CONES}
(\code{feedback\_struct.h:69}) and $q_{\rm e}$ is \code{STAR\_GAS\_WEIGHT\_EDGE}
(\code{src/feedback/GEAR/star\_gas\_weights.h:34}). Eq.~\eqref{eq:packet-tent} is
\code{star\_gas\_tent\_weight} (\code{star\_gas\_weights.h:54-58}),
Eq.~\eqref{eq:packet-floor} is \code{star\_gas\_cone\_floor} (:71-76, with $f_0$ the
\code{STAR\_GAS\_CONE\_FLOOR\_FRACTION} of :38), Eq.~\eqref{eq:packet-emit-share}
is \code{star\_gas\_global\_share} (\code{star\_gas\_weights.h:125-134}), which takes
$D$ from \code{star\_gas\_escape\_fraction} (:103-109), the sum over the cells of
Eq.~\eqref{eq:packet-escape}, \code{star\_gas\_cone\_escape} (:88-92). It is called in
\code{radiation\_packet\_star\_emit} (\code{radiation\_packet.h:513-515}) with the 12 cell sums
and the floor (\code{radiation\_packet\_cone\_floor}, :436). The energy and the photon number
follow from the fraction, then the delivery. The star zeroes its cone
sums before the density loop (\code{radiation\_packet\_init\_star\_cones}, called
at \code{src/feedback/GEAR\_thermal/feedback.c:580}).

\subsection{Redshift and band-edge transfer}
\label{sec:packet-transfer}

In an expanding box a share loses energy to the expansion along its path. The
loss is taken once, when the share is booked, and before the dust acts
(Section~\ref{sec:packet-absorption}). The weights are those of the M1 system
(Section~\ref{sec:isrf-eqdisc-lambda}). A band edge is fixed in physical
energy, so a band loses photons through its lower edge: the energy of band $b$
falls with the band-edge weight $\lambda_E(b)\ge1$, and the part that crosses
the $11.2$\,eV edge from LW enters PE. With $\ell_{\rm phys}$ the physical path of the share
(Eq.~\eqref{eq:packet-path} times $a$), $H$ the Hubble rate and $c$ the speed of light, the redshift
depth is
\begin{equation}
  \zeta = \frac{H\,\ell_{\rm phys}}{c}.
  \label{eq:packet-zeta}
\end{equation}
The energy of the share in LW and PE, before ($e$) and after ($e'$), is
\begin{align}
  e'_{\rm LW} &= e_{\rm LW}\,e^{-\lambda_E({\rm LW})\,\zeta}, \nonumber\\
  e'_{\rm PE} &= e_{\rm PE}\,e^{-\lambda_E({\rm PE})\,\zeta}
    + \frac{\lambda_E({\rm LW})-1}{\lambda_E({\rm LW})}\,
      \big(e_{\rm LW}-e'_{\rm LW}\big).
  \label{eq:packet-redshift-E}
\end{align}
The photon number $n$ of LW (stored as $N E_{\rm lo}$) falls with its own weight $\lambda_N$, and the photons
that it loses enter PE in units of the PE lower edge. The photon number of PE
falls at the rate with which its own photons cross $6$\,eV, with
$\nu=(\lambda_E({\rm PE})-1)\,e_{\rm PE}/n_{\rm PE}=(\lambda_E({\rm PE})-1)\,\bar E_{\rm PE}/E_{{\rm lo},{\rm PE}}$,
$\bar E_{\rm PE}$ being the mean photon energy of the share:
\begin{align}
  n'_{\rm LW} &= n_{\rm LW}\,e^{-\lambda_N\,\zeta}, \nonumber\\
  n'_{\rm PE} &= n_{\rm PE}\,e^{-\nu\,\zeta}
    + \frac{E_{{\rm lo},{\rm PE}}}{E_{{\rm lo},{\rm LW}}}\,
      \big(n_{\rm LW}-n'_{\rm LW}\big).
  \label{eq:packet-redshift-N}
\end{align}
No table gives $\nu$, because the M1 system has no PE photon moment. The
photon gain of PE and its energy gain describe the same photons only if the
mean photon energy of the share is the table mean from which the weights come.
Only the LW photon number reaches Grackle (Section~\ref{sec:packet-field}), so
the PE photon number is a ledger quantity. In every band, the amount removed,
net of the gain from the other band, is added to the redshift ledger $Y$ of the
receiver (Section~\ref{sec:packet-ledger}). $Y_{\rm PE}$ can be negative. At
fixed $a$ and $H$, and when each merge receives a single share, $\zeta$ adds along
a chain of hops as the path does. One long hop and many short hops then give the
same LW energy and LW photon number, and the same PE amounts to first order in
$\zeta$. $H$ and $a$ are the values at the time of the booking, for a hop and for
a star. With $H=0$ in a
static box, there is no redshift and no transfer.

The M1 system takes one combined depth for the dust and the expansion, and
splits the LW loss by the ratio of the two rates (Eq.~\eqref{eq:isrfd-transfer-fedge}).
The packet transport applies the redshift first and the dust second. The two agree in
the energy that LW transmits, $E_{\rm LW}e^{-(\tau+\lambda_E({\rm LW})\zeta)}$, and
in the PE gain when $\kappa=0$. With dust, the PE gain and the split of the loss between the
redshift ledger and the absorbed ledger differ by a term of order $\zeta\tau E$.

\paragraph{Implementation.} The constants are set at start-up by
\code{radiation\_packet\_startup} (\code{radiation\_packet.c:93-99}, called at
\code{src/runner\_radiation\_feedback.c:240}) through
\code{radiation\_packet\_set\_redshift} (:77-86), from the band-edge weights
\code{band\_edge\_weight\_pe}, \code{band\_edge\_weight\_lw} and
\code{band\_edge\_photon\_weight\_lw}
(\code{src/feedback/GEAR\_thermal/feedback\_properties.h:225-227}).
Eqs.~\eqref{eq:packet-zeta}--\eqref{eq:packet-redshift-N} are
\code{radiation\_packet\_redshift\_share} (\code{radiation\_packet.c:122-168}):
$\zeta$ at :134, $\nu$ at :143-144, the PE photon loss at :145, the losses at :147-151, the gains at
:154-157, the update at :159-162 and the returned ledger terms at :164-167.
\code{radiation\_packet\_deliver} calls it when $H>0$ and adds the terms to
\code{redshifted\_energy} and \code{redshifted\_photons}
(\code{radiation\_packet.h:225-239}; fields at \code{feedback\_struct.h:274-278}).
The path is multiplied by $a$ by the callers (\code{radiation\_packet.h:340} for
a hop, :522 for a star), and $H$ comes from
\code{runner\_iact\_isrf\_packet\_hop} (\code{radiation\_propagation\_iact.h:1043-1044},
:1075) and from \code{radiation\_iact.h:275-277}. The M1 transfer is at
\code{radiation\_isrf.c:639-671}; the packet transport returns before it
(\code{radiation\_isrf.c:610-614}).

\subsection{Dust absorption}
\label{sec:packet-absorption}

A share is absorbed by dust along the distance that its light travelled
in the receiver. In an expanding box the redshift of
Section~\ref{sec:packet-transfer} acts first, and the dust acts on what remains. The dust uses the linear dust absorption rate $\kappa_b$ of the
receiver in band $b$, the same $\kappa_b$ that the M1 system uses
(Section~\ref{sec:isrf-eqdisc-continuum}). The optical depth, the absorbed
fraction and the transmitted amounts are
\begin{align}
  \tau_b &= \kappa_b\,\ell_{\rm phys} = \kappa_b\,a\,\Delta\ell, \label{eq:packet-tau}\\
  \alpha_b &=
  \begin{cases}
    1-e^{-\tau_b}, & \tau_b<80,\\
    1, & \tau_b\ge 80,
  \end{cases} \label{eq:packet-alpha}\\
  e_{b}^{\rm tr} &= (1-\alpha_b)\,e_b, \qquad
  n_{b}^{\rm tr} = (1-\alpha_b)\,n_b, \label{eq:packet-transmit}
\end{align}
with $e_b$ and $n_b$ the energy and the photon number of the share after the
redshift and before the absorption. Energy and photon number lose the same fraction, so the mean photon
energy of a band is not changed by absorption (Eq.~\eqref{eq:packet-meanE}).
If $n^{\rm tr}_b<10^{-35}$ (internal units of $N E_{\rm lo}$), the whole share of
the band is booked as absorbed and nothing is transmitted. This keeps the photon number
away from the float minimum, where it would flush to zero before the energy.

The path length is
\begin{equation}
  \Delta\ell =
  \begin{cases}
    \max(0,\,d_r-d_s), & \text{hop from a particle},\\
    r_{\star j}, & \text{emission by a star},
  \end{cases}
  \label{eq:packet-path}
\end{equation}
where $d_s=d_{i,k}$ is the distance from the origin to the sender, and
$d_r=|\boldsymbol{o}_j|$ is the distance from the origin to the receiver,
$\boldsymbol{o}_j$ being the offset of Eq.~\eqref{eq:packet-offset}. For a
star, $d_s=0$ and $d_r=r_{\star j}$. For $\mu_{ij,k}>0$, $d_r^2=d_s^2+2d_s\,\boldsymbol{\Delta x}_{ij}\!\cdot\!\boldsymbol{n}_{i,k}+r_{ij}^2>d_s^2$, so the clamp acts only on rounding. When each merge
receives a single share, the merge keeps the origin of that share, and the path lengths
add up to the distance from the origin to the last receiver. The transmitted amounts are booked into the inbox of
the receiver slot, or into its reservoir (Section~\ref{sec:packet-reservoir}), and the
absorbed amounts into the absorbed ledger of the receiver
(Section~\ref{sec:packet-ledger}). The absorbed energy is a diagnostic: it is
not added to the gas, which sees the light only through the field of
Section~\ref{sec:packet-field}.

\paragraph{Implementation.} All of it is in
\code{radiation\_packet\_deliver} (\code{radiation\_packet.h:196-287}):
the redshift at :225-239, before the dust;
Eq.~\eqref{eq:packet-tau} at :246, Eq.~\eqref{eq:packet-alpha} at :247-249
(\code{RADIATION\_PACKET\_MAX\_OPTICAL\_DEPTH} is :55, and the exponential is
\code{expm1}, in double), Eq.~\eqref{eq:packet-transmit} at :251-254, the
photon cutoff at :255-260 (\code{RADIATION\_PACKET\_MIN\_PHOTONS}, :60), and
the ledger at :269-270. The rate $\kappa_b$ is
\code{isrf\_operator[b].kappa} of the receiver
(\code{feedback\_struct.h:168-171}), set once per step by
\code{radiation\_snapshot\_part\_propagation} (\code{radiation\_isrf.c:216-223}), as in the M1 system. The hop
path, Eq.~\eqref{eq:packet-path} first line, is at
\code{radiation\_packet.h:328-331}, with $a$ applied at the call, :340; the star path is
:521-522 (call at :522).

\subsection{Particles on different time bins}
\label{sec:packet-reservoir}

A share goes into the inbox of the receiver slot, from which it joins the
held packet at the next active step of the receiver. This holds if the
sender's time bin is not coarser than the receiver's. If the sender's bin is
coarser (its step is longer), the share goes into a reservoir of the slot
instead, with the same origin sum, and the time $T_{\rm end}$, the time of the booking plus one step of the source's current time bin, is recorded
(the latest over all shares of the slot). The reservoir is released in equal
parts over the sender's step. At each merge of the receiver, at time $t$ and after a step of length $\Delta t$ (the step that ends now, which began at $t-\Delta t$), the fraction
\begin{equation}
  f_{{\rm rel},k} = \begin{cases}
    \Delta t/(T_{\rm end}-t+\Delta t), & T_{\rm end}>t,\\
    1, & \text{otherwise}
  \end{cases}
  \label{eq:packet-release}
\end{equation}
of the reservoir, measured from the start of the step that ends now (under cosmology $\Delta t$ and $T_{\rm end}-t+\Delta t$ are proper times), of its energy, its photon number and its origin sum alike, joins the arrivals
of the merge (Eq.~\eqref{eq:packet-merge-E}). With equal receiver steps the parts
are equal: a share booked at the start of a receiver step, with $n\Delta t$ to run, is released
in $n$ merges of $1/n$ each, and the last merge releases the rest. If the release of the fraction $f_{{\rm rel},k}$ would leave the photon
number of a band below $10^{-35}$, the whole reservoir is released. A star
sends to the reservoir of a receiver whose bin is finer than its own, with
$T_{\rm end}$ equal to the time of the emission plus one step of its time bin. This reservoir is a transport buffer for
unequal steps, and it is not the dose reservoir of the M1 injection
(Section~\ref{sec:isrfd-reservoir}).

\paragraph{Implementation.} The choice between the inbox and the reservoir is
the argument \code{to\_reservoir} of \code{radiation\_packet\_deliver}
(\code{radiation\_packet.h:196-287}); it is
\code{s->time\_bin > r->time\_bin} at :340 for a hop, and
\code{si->time\_bin > pj->time\_bin} at :523 for a star, with
\code{source\_step\_end\_ti} (:201, :277-278) the end of the sender's step
(\code{step\_end\_ti}, set at \code{radiation\_packet.c:718} as \code{ti\_current + ti\_step}, and
\code{ti\_current + ti\_step} for a star, \code{radiation\_packet.h:523}). Eq.~\eqref{eq:packet-release} is
\code{radiation\_packet\_merge\_slot} (\code{radiation\_packet.c:605-617}); the
cutoff rule is at :619-627, and the release of the energy, the photon number and the
origin sum at :629-665.

\subsection{The field that the gas sees}
\label{sec:packet-field}

The held packet of a slot is the light that crosses the particle in one step
over the advance $l_{k}$ of Eq.~\eqref{eq:packet-advance}. The specific energy
of band $b$ that the particle sees is the sum over its slots,
\begin{equation}
  u_{b} = \sum_{k:\,S_k>0} \frac{E_{k,b}}{m}\,\frac{a\,l_k}{c\,\Delta t},
  \label{eq:packet-field}
\end{equation}
with $m$ the particle mass, $\Delta t$ the physical length of the step that ends now and $c$
the speed of light; the photon number $N_{k,b}$ gives the field of the photon
number in the same way. The factor $a\,l_k/(c\Delta t)$ is the hop speed over $c$ (Section~\ref{sec:packet-frame}). It is
set to zero for $\Delta t\le0$ or $m\le0$. The field can be averaged in time with an
exponential running mean of time constant $\tau_{\rm m}$
(\code{ISRF\_packet\_read\_mean\_time\_Myr}, default 0, i.e.\ off),
\begin{equation}
  \bar u \leftarrow \bar u + (u-\bar u)\,\frac{\Delta t}{\tau_{\rm m}}, \qquad
  \bar u \leftarrow u \ \text{if}\ \tau_{\rm m}\le0\ \text{or}\ \Delta t\ge\tau_{\rm m}.
  \label{eq:packet-runmean}
\end{equation}
A step with $\Delta t\le0$ keeps the old mean.
Only the field is averaged and never the packet, so the ledger is not
changed. Then, in the force loop, each particle $i$ takes the kernel mean of the
fields of its neighbours and of itself,
\begin{equation}
  \langle u_b\rangle_i =
  \frac{\sum_j m_j\,w(r_{ij}/h_i)\,\bar u_{b,j} + m_i\,w(0)\,\bar u_{b,i}}
       {\sum_j m_j\,w(r_{ij}/h_i) + m_i\,w(0)},
  \label{eq:packet-kernelmean}
\end{equation}
over the neighbours with $r_{ij}<\gamma h_i$, with $w(u)$ the kernel of the build
without the factor $h^{-3}$. The kernel mean $\langle u_b\rangle_i$ is written into the
band moments that Grackle reads: PE to PE, LW to LW and the photon number
of LW to the LW photon-number moment, in place of any value that was there. The photon number of PE is not coupled.

\paragraph{Implementation.} Eq.~\eqref{eq:packet-field} is
\code{radiation\_packet\_end\_gradient} (\code{radiation\_packet.c:742-791}, the
sum over slots at :762-773, the factor at :767); it is called from
\code{radiation\_end\_gradient\_propagation} (\code{radiation\_isrf.c:924-930}).
Eq.~\eqref{eq:packet-runmean} is \code{radiation\_packet\_read\_mean}
(\code{radiation\_packet.h:113-118}), applied at
\code{radiation\_packet.c:785-790}. The sums of
Eq.~\eqref{eq:packet-kernelmean} are \code{radiation\_packet\_gather\_field}
(\code{radiation\_packet.h:369-386}, called at
\code{radiation\_propagation\_iact.h:1045-1046}, :1072), and the self term and the
division are \code{radiation\_packet\_end\_force} (\code{radiation\_packet.c:799-818}),
called at \code{radiation\_isrf.c:611}. The write into the moments is
\code{radiation\_packet\_couple\_field} (\code{radiation\_packet.c:831-839}, called at
\code{radiation\_isrf.c:612}). The packet transport sets no timestep limit from a
propagation speed (\code{radiation\_isrf.c:316-318}).

\subsection{The frame of reference of the light}
\label{sec:packet-frame}

The packet transport works in the gas frame. Every offset is measured from the
particle that carries the packet: the stored origin of a slot is a distance and
a direction from its carrier (Section~\ref{sec:packet-merge}), and the offset
of a share is relative to its receiver (Eq.~\eqref{eq:packet-offset}). No
position in the box frame enters the light. The distance and the direction of a
held packet are set at the merge, which runs at the drift, and the hop of the
same integer time reads them before the particle has moved, so they are never
stale. The offsets of the inbox and of the reservoir are different. They are
measured from the receiver at the time of booking and stay as they are until
the next merge of the receiver, over any drift of the receiver in between. This
is the gas-frame choice.

The M1 system makes the same choice. Its equation is written along the motion
of the gas, so the advection is supplied by the particle itself, and the
remaining $v/c$ terms are dropped one by one
(Sections~\ref{sec:fuv-lagrangian} and \ref{sec:fuv-assumptions}).
GEAR-RT keeps the radiation field of a particle in a static frame and corrects
it at each drift by extrapolation along the gradient.

Real light moves at $c$ in every frame. The gas frame and the laboratory frame
differ by terms of order $v/c$, which both schemes drop. The choice matters
because a hop has the speed $v_{\rm hop}=a\,l/\Delta t$, with $l$ the
advance of Eq.~\eqref{eq:packet-advance} and $\Delta t$ the step of the sender.
This speed is usually far below $c$, but the code does not bound it by $c$
(Section~\ref{sec:packet-hop}).
In a laboratory (comoving) frame, a source that moves with its gas at the bulk
speed $V$ would leave its light behind, by about $V/v_{\rm hop}$ hop lengths per
hop, and a uniform boost of the whole box would change the result. In the gas
frame, only the velocity of the source relative to the absorbing gas enters,
and it is small for a star that moves with its natal cloud. The Hubble flow is
a coordinate effect that comoving offsets absorb exactly.

Per hop, the terms that the transport drops are of order $\Delta v/c$, with
$\Delta v$ the relative velocity of the sender and the receiver: $10^{-5}$ to
$10^{-4}$ for $\Delta v=3$ to $30\,{\rm km\,s^{-1}}$. The redshift that the
transport keeps is the Doppler shift of the Hubble flow over the same hop. For a
$7$\,pc hop it is $\lambda\zeta$, with $\lambda$ the band-edge weight and $\zeta=1.6\times10^{-9}$ at $z=0$ and
$2.8\times10^{-8}$ at $z=9$ ($H_0=70\,{\rm km\,s^{-1}\,Mpc^{-1}}$,
$\Omega_{\rm m}=0.3$). The ratio of the two is the velocity gradient over $H$.
Inside a galaxy this ratio is $10^{2}$ to $10^{3}$: $10\,{\rm km\,s^{-1}}$ over
$100$\,pc is $3\times10^{-15}\,{\rm s^{-1}}$, against $H_0=2.3\times10^{-18}\,{\rm s^{-1}}$ and
$H(z=9)=3.9\times10^{-17}\,{\rm s^{-1}}$. In diffuse intergalactic gas, where the
peculiar gradient is of order $H$, the two are comparable.

Two other choices are possible. A laboratory-frame cone would subtract the drift
displacement from the stored inbox and reservoir offsets at each drift, for 0
to 24\,B per gas particle. A correction for the shear of the flow would also
move the offsets with the velocity gradient, for 192 to 384\,B per gas
particle.

\paragraph{Implementation.} The merged origin is set in
\code{radiation\_packet\_merge\_slot} (\code{radiation\_packet.c:669-691}) at the
drift, \code{radiation\_snapshot\_part\_propagation} (\code{radiation\_isrf.c:257-266}). The hop
reads it in \code{radiation\_packet\_book\_share} (\code{radiation\_packet.h:325-327}).
An empty slot keeps its old origin (\code{radiation\_packet.c:667}), but the
gradient and the hop skip empty slots (\code{radiation\_packet.h:159, :311}). The
inbox and reservoir sums are booked at \code{radiation\_packet.h:274-282}. In
the M1 chapter the advection is at \code{02\_fuv\_isrf.tex:457-468} (Eq.~\eqref{eq:fuv-lagr-line2}) and the
gas-frame argument at :675-682; the dropped terms are
at :302-351. GEAR-RT is described in
\code{theory/RadiativeTransfer/GEARRT/3-rt-numerical-method.tex}, lines 225-228
(static frame) and 750-771 (Section ``Dealing With Particle Drifts'').

\subsection{Ledgers}
\label{sec:packet-ledger}

The light is accounted for per band, in energy and in photon number ($N E_{\rm lo}$).
Every gas particle carries a cumulative absorbed ledger $A$ (Section~\ref{sec:packet-absorption}),
a cumulative redshift ledger $Y$ (Section~\ref{sec:packet-transfer}; zero in a static box)
and a cumulative escaped ledger $X$: the held light of a slot with no
downstream neighbour (Section~\ref{sec:packet-hop}) is added to $X$ and is set to zero.
Every star carries the energy emitted, $E_{\rm em}=L_b\Delta t_\star$ summed over
its steps, and the energy escaped, $X_\star=L_b\Delta t_\star\sum_c\varepsilon_c$
(Eq.~\eqref{eq:packet-escape}), together with the photon numbers $f_b$ times these.
When a gas particle becomes a star, or is swallowed by a sink or becomes one, the light it still
holds in flight, which is the held, the inbox and the reservoir amounts of all slots, is
lost with it, and is booked as removed energy $D$ in the ledger of the
star or the sink. The absorbed, the redshift and the escaped ledgers of the gas are carried over in
the same ledger as inherited amounts $I_{\rm inh}$. When a sink swallows another, the
survivor adds the ledger of the swallowed sink to its own. With $E_{\rm fl}$ the light in flight,
the amounts of one band balance as
\begin{equation}
  \sum_\star E_{\rm em} =
  \sum_{\rm gas}\big(A+X+Y+E_{\rm fl}\big) + \sum_{\star,\,{\rm sinks}}\big(X_\star + D + I_{\rm inh}\big),
  \label{eq:packet-balance}
\end{equation}
and likewise for the photon number. In each band, $Y$ is booked net of the
gain from the other band, so $Y_{\rm PE}$ can be negative, and the balance of each band is exact. Each step of the transport moves light
between the terms without creating or destroying it (the hop shares sum to
one, Eq.~\eqref{eq:packet-share}; the merge adds; the redshift moves the lost
part to $Y$; the absorption moves the absorbed part to $A$). The cone fractions of a star sum to one only if its
density loop and its feedback loop see the same receivers and the same $h_\star$.
The balance holds up to the rounding of float arithmetic in a periodic box,
or while no gas particle leaves the box. Gas that leaves a non-periodic box is
removed with its light in flight and its ledgers, and nothing books them.
The ledgers are in double, because small shares must
not be lost against a large total. A stored share is a float, with a relative
rounding of at most $2^{-24}$, about $6\times10^{-8}$. The relative loss of a
share in one hop is $\lambda\zeta$, with $\lambda$ the band-edge weight. For a $7$\,pc hop with the shipped weights it
is about $10^{-8}$ in LW at $z=0$, and $1.8\times10^{-7}$ in LW and $6\times10^{-8}$ in PE at
$z=9$. Below $z$ of about 4 the LW loss of one hop is below the rounding, so the
stored shares carry it only on average. At higher redshift they resolve it. The
double ledger $Y$ books it in both cases. A bar on the energy balance must lie
above that rounding. The test beams of the code, which add light
without a star, are in no ledger and are not part of Eq.~\eqref{eq:packet-balance}.

\paragraph{Implementation.} The per-band ledger of a gas particle is
\code{feedback\_isrf\_packet\_band\_data} (\code{feedback\_struct.h:242-279}:
\code{absorbed\_energy}, \code{absorbed\_photons}, \code{escaped\_energy},
\code{escaped\_photons}, \code{redshifted\_energy}, \code{redshifted\_photons}); the ledger of a star or a sink is
\code{feedback\_isrf\_packet\_removed\_data}
(\code{feedback\_struct.h:378-391}), and the emission and escape of a star are
\code{emitted\_energy}, \code{emitted\_photons}, \code{escaped\_energy} and
\code{escaped\_photons} in the \code{packet} struct of the star
(\code{feedback\_struct.h:808-852}). The emission is booked by
\code{radiation\_packet\_star\_book\_emission} (\code{radiation\_packet.c:852-882}:
emitted at :858-865, escaped at :874-881, $\sum_c\varepsilon_c$ at
:866-872), called at \code{src/feedback/GEAR\_thermal/feedback.c:649}. The
removal is \code{radiation\_packet\_book\_removed\_gas}
(\code{radiation\_packet.h:543-562}, the inherited sum with the redshift ledger at :557-560), reached from
\code{feedback\_book\_removed\_gas\_to\_spart},
\code{feedback\_book\_removed\_gas\_to\_sink} and
\code{feedback\_book\_sink\_merger} (\code{src/feedback/GEAR/feedback\_common.c:1785-1810}).
The escape of a slot with no downstream neighbour is at
\code{radiation\_packet.c:775-782}.
